\documentclass[12pt,a4paper]{article}
\usepackage[T1]{fontenc}
\usepackage{amsmath}

\begin{document}
	After the first permutation we have the structure
	$$
	\mathcal{A} =
	\begin{pmatrix}
		.A & B \\
		C & D
	\end{pmatrix}
	$$
	or
	$$
	\begin{pmatrix}
		.A_{1} & 0 & \cdots & 0 & B_{1} \\
		0 & .A_{2}  &  \cdots & 0 & B_{2} \\
		\vdots & \vdots &  \vdots  &  \vdots  &  \vdots \\
		0 & .0  &  \cdots & A_{n} & B_{n} \\
		C_1 & C_2 & \cdots  & C_{n} & D
	\end{pmatrix}
	$$
	We want now solve a linear system
	$$
		\mathcal{A} 	x = b
	$$
	with the partition 
	$$
	x = 
	\begin{pmatrix}
		x_1 \\
		x_2
	\end{pmatrix}
	$$ 
	and
	$$
	x_1=
	\begin{pmatrix}
		x_{1i} \\
		\vdots \\
		x_{1n}
	\end{pmatrix}
	$$ 
	and respectively for $b$.
	
	Computation of the Schur complement matrx 
	\begin{align*}
		S& = D - \sum_i C_iA_i^{-1}B_i
	\end{align*}
	
	\begin{align*}
		x_{1i}  & = A_i^{-1}b_{1i} \quad i=1,\ldots,n\\
		x_2 & = S^{-1}(b_2-\sum_i C_ix_{1i})\\
		x_{1i} & = A_i^{-1}(b_{1i}-B_ix_2), \quad i=1,\ldots,n
	\end{align*}
	
	The matrices $A_i$ are also in block structured where the upper left block is the identity matrix.
	

\end{document}